% ============================================================================
% UTILS LaTeX - Snippet utili per la scrittura tesi
% Nota: questo file e' una raccolta di esempi da copiare/incollare.
% ============================================================================

% -------------------------
% Struttura documento
% -------------------------
\chapter{Titolo Capitolo}
\label{chap:capitolo}

\section{Titolo Sezione}
\label{sec:sezione}

\subsection{Titolo Sottosezione}
\label{subsec:sottosezione}

\subsubsection{Titolo Sottosottosezione}
\label{subsubsec:tema}

% -------------------------
% Riferimenti incrociati
% -------------------------
Vedi il Capitolo~\ref{chap:capitolo}, la Sezione~\ref{sec:sezione}
e la Figura~\ref{fig:esempio}.

% -------------------------
% Liste
% -------------------------
\begin{enumerate}
\item primo oggetto
\item secondo oggetto
\item terzo oggetto
\item quarto oggetto
\end{enumerate}

\begin{itemize}
\item primo oggetto
\item secondo oggetto
\end{itemize}

\begin{description}
\item[OGGETTO1] prima descrizione.
\item[OGGETTO2] seconda descrizione.
\item[OGGETTO3] terza descrizione.
\end{description}

% -------------------------
% Figure
% -------------------------
\begin{figure}[htbp]
\centering
\includegraphics[width=0.60\textwidth]{images/nome_immagine.png}
\caption[Legenda breve per indice]{Legenda completa sotto la figura.}
\label{fig:esempio}
\end{figure}

% Due immagini affiancate (richiede pacchetto subcaption)
% \usepackage{subcaption}
\begin{figure}[htbp]
\centering
\begin{subfigure}[b]{0.48\textwidth}
\centering
\includegraphics[width=\textwidth]{images/fig_a.png}
\caption{Caso A}
\label{fig:caso-a}
\end{subfigure}
\hfill
\begin{subfigure}[b]{0.48\textwidth}
\centering
\includegraphics[width=\textwidth]{images/fig_b.png}
\caption{Caso B}
\label{fig:caso-b}
\end{subfigure}
\caption{Confronto tra due casi.}
\label{fig:confronto-casi}
\end{figure}

% -------------------------
% Tabelle
% -------------------------
\begin{table}[htbp]
\centering
\begin{tabular}{r|c|c}
\hline \hline
$(1,1)$ & $(1,2)$ & $(1,3)$\\
\hline
$(2,1)$ & $(2,2)$ & $(2,3)$\\
\hline
$(3,1)$ & $(3,2)$ & $(3,3)$\\
\hline \hline
\end{tabular}
\caption[Legenda breve tabella]{Legenda completa tabella.}
\label{tab:uno}
\end{table}

% Variante consigliata (richiede booktabs)
% \usepackage{booktabs}
% \begin{tabular}{lcc}
% \toprule
% Colonna A & Colonna B & Colonna C \\
% \midrule
% Valore 1 & 10 & 20 \\
% Valore 2 & 12 & 18 \\
% \bottomrule
% \end{tabular}

% -------------------------
% Formule
% -------------------------
Formula inline: $E = mc^2$.

\begin{equation}
f(x) = ax^2 + bx + c
\label{eq:parabola}
\end{equation}

\begin{align}
y &= mx + q \\
z &= \alpha x + \beta
\end{align}

% -------------------------
% Citazioni e bibliografia
% -------------------------
Riferimento singolo: \cite{K1}.
Riferimento multiplo: \cite{K1,K2,K3}.

% -------------------------
% URL, note, enfasi
% -------------------------
% (richiede hyperref)
Link: \url{https://www.example.com}
Link con testo: \href{https://www.example.com}{Sito di esempio}

Testo in \textbf{grassetto}, \textit{corsivo} e \textsc{maiuscoletto}.
Nota a pie' di pagina\footnote{Questa e' una nota di esempio.}.

% -------------------------
% Codice/listati
% -------------------------
\begin{verbatim}
In questo ambiente posso scrivere come voglio,
incluso il simbolo % senza commentare.
\end{verbatim}

% Alternativa migliore per codice: listings
% \usepackage{listings}
% \begin{lstlisting}[language=Python, caption={Esempio Python}, label={lst:py-esempio}]
% def hello(name):
%     return f"Ciao {name}"
% \end{lstlisting}

% -------------------------
% Comandi personalizzati
% -------------------------
% Da mettere nel preambolo del file principale:
% \newcommand{\pmbok}{PMBOK~8}
% \newcommand{\prince}{PRINCE2}
% \newcommand{\agile}{Agile}

% Uso nel testo:
% In questo capitolo confrontiamo \pmbok{} e \prince{}.